\documentclass[11pt,a4paper]{article}
\newcommand{\labonaamwaarde}{lab 4}
\usepackage[english]{syllabus}

\begin{document}

\syllabustitel{lab 4 -- \emph{Data Processing II}}{Combine, summarize, reshape}

% ---------------------------------------------------------------------------
% Shared by the syllabi of labs 3, 4 and 5. Edit it here, not in the lab files.
\section{Overview}

In labs 3, 4 and 5 you are the new data analyst of a (fictional) youth
football and athletics club. The head coach gives you three export files
from the club's systems. They are messy on purpose, the way real exports
are: inconsistent labels, sprint times in milliseconds, missing values,
an athlete who appears twice and a column with a comma as the decimal
separator.

Each session builds on the previous one:

\begin{description}[leftmargin=1.6cm, style=nextline]
  \item[Lab 3] \textbf{From file to tidy table.} Load, explore, filter,
    add columns and clean until you can trust the data.
  \item[Lab 4] \textbf{Combine, summarize, reshape.} Summarize groups,
    combine tables, switch between a wide and a long shape, and work with
    a heart-rate series per second.
  \item[Lab 5] \textbf{Tell the story with plotly.} Choose the right chart
    for a question, make honest charts with \code{plotly.express} and
    combine them into a small interactive dashboard.
\end{description}


% ---------------------------------------------------------------------------
\section{Learning outcomes}

After lab 4 you can:

\begin{enumerate}[label=\textbf{LO\arabic*}, leftmargin=1.4cm]
  \item summarize data per group with \code{groupby} and \code{agg}, and
        combine tables with \code{merge} without double counting;
  \item reshape data between a long and a wide shape with \code{melt}
        and \code{pivot\_table};
  \item compute basic features of a time series: heart-rate zones,
        rolling averages and resampling.
\end{enumerate}

% ---------------------------------------------------------------------------
\section{Prerequisites}

Lab 3 (\emph{Data Processing I}): loading a CSV file with \code{pandas},
selecting and filtering rows and columns, computing new columns and
cleaning the data. Lab 4 starts with the cleaning pipeline of lab 3, in
compact form.

% ---------------------------------------------------------------------------
% Shared by the syllabi of labs 3, 4 and 5. Edit it here, not in the lab files.
\section{Software and materials}

Every lab is a \textbf{marimo notebook} in the \code{labs/} folder of the
course repository. You can open it in two ways:

\begin{itemize}
    \item \textbf{Online:} go to \url{https://marimo.app}, choose
        \emph{New $\rightarrow$ Open from URL\dots} and paste the raw URL
        of the notebook on GitHub. The data is then read from GitHub, so
        you don't have to install anything.
    \item \textbf{Locally:} install the packages from
        \code{requirements.txt} and open the notebook with marimo. The
        data is then read from the \code{data/} folder.
\end{itemize}

Three things to know about marimo:

\begin{itemize}
  \item Notebooks are \textbf{reactive}: when you change a cell, every
        cell that depends on it runs again, in the order of their
        dependencies.
  \item \textbf{Two cells cannot define the same variable.} Give every
        step its own clear name, such as \code{tests\_raw} and
        \code{tests\_clean}, or keep a sequence of steps together in one
        cell.
  \item Every exercise contains a \code{...} placeholder. A hidden cell
        below it checks your answer: it turns \textbf{green} when the
        answer is right, and otherwise gives a general hint about what is
        still wrong, without giving the answer away.
\end{itemize}


% ---------------------------------------------------------------------------
% Shared by the syllabi of labs 3, 4 and 5. Edit it here, not in the lab files.
\section{The dataset}

\noindent
\rowcolors{2}{tint}{white}
\begin{tabularx}{\textwidth}{@{}L{3.1cm}L{3.2cm}Y@{}}
\toprule
\textbf{File} & \textbf{Size} & \textbf{Content and known problems} \\
\midrule
\code{athletes.csv}   & 43 rows, 7 columns   & ID, name, birthdate, gender, position, height, weight. 13 spellings for 4 positions, weight as text with a decimal comma (\code{"57,5"}), one athlete entered twice with two IDs, a few missing values. \\
\code{tests.csv}      & 114 rows, 8 columns  & 10\,m and 30\,m sprint, countermovement jump height and Yo-Yo level, at 3 test moments. Some sprint times are stored in milliseconds; a few missing values. \\
\code{hr\_session.csv} & 12,000 rows, 5 columns & Heart rate and speed per second for 5 athletes during one 40-minute training session. \\
\bottomrule
\end{tabularx}
\rowcolors{1}{}{}


% ---------------------------------------------------------------------------
\needspace{18\baselineskip}
\section{Schedule}

\textbf{Goal:} aggregate data with \code{groupby}, combine tables with
\code{merge}, reshape with \code{melt} and \code{pivot}, and compute
basic features of a time series.

\begin{schedule}
13:00--13:15 & Recap: repair a broken cell & Fix a small cleaning pipeline with two bugs. \\
13:15--14:10 & \code{groupby} + \code{agg} & Statistics per group, age groups with \code{pd.cut}, several named aggregations at once. \\
14:10--14:55 & \code{merge} & Inner and left joins; why a duplicated athlete is counted twice, and why you remove duplicates before merging. \\
14:55--15:10 & Break & \\
15:10--15:50 & \code{melt} and \code{pivot} & Long and wide shape; progress between two test moments. \\
15:50--16:50 & Time series & Heart-rate zones as \% of $HR_{max}$ (220 $-$ age), time per zone, \code{rolling}, \code{resample}. \\
16:50--17:30 & Mini-case in pairs & A training report: one row per athlete. \\
\end{schedule}

\textbf{Exercises:} (0) recap: repair the broken cell; (1) group by
gender and age group; (2) what did removing the duplicates cost?;
(3) progress in jump height; (4) minutes in zone 4--5 per athlete;
\textbf{mini-case:} training report with the minutes in zone 4--5, the
maximum speed and the distance covered per athlete.

% ---------------------------------------------------------------------------
\section{Working method}

\begin{itemize}
  \item Each block starts with a short demo in the notebook, followed by
        exercises you solve in the notebook itself.
  \item Use the automatic checks as feedback, not as the goal: when a
        check turns green, make sure you can explain \emph{why} your code
        works.
  \item You do the mini-case at the end in pairs.
\end{itemize}

% ---------------------------------------------------------------------------
\section{Preparation}

\begin{itemize}
  \item Check before the session that you can open the notebook, locally
        or on marimo.app.
  \item Make sure you can explain every step of the cleaning pipeline at
        the end of lab 3, because lab 4 starts from it.
  \item Read the course text for lab 4 next to the notebook.
\end{itemize}

\end{document}
